
\documentclass[11pt]{article}

\usepackage{pgf}
\usepackage{tikz}
\usetikzlibrary{arrows,automata}
\usepackage{qtree}
\usepackage{url}
\usepackage[margin=2.5cm]{geometry}
\usepackage{amsmath}
\usepackage{amssymb}
\setlength{\parindent}{0pt}

\begin{document}

\rule{\textwidth}{1pt}
\begin{center}
\textbf{Taaltheorie en Taalverwerking -- Homework Submission}
\end{center}
\rule{\textwidth}{0.5pt}\\

\textbf{Week:} ...

\textbf{Group:} ...

\textbf{Student names and IDs:} ...

\textbf{Date:} ...

\rule{\textwidth}{1pt}

\begin{enumerate}

\item
\begin{enumerate}
\item

\begin{tabular}{l|c|c|c|c|c}
Rijksmuseum & doc1 & doc2 & doc3 & doc4 & q = what is the rijksmuseum in Amsterdam, Netherlands \\
\hline
Amsterdam & 1 & 1 &  & 1 & 1\\
Netherlands & 1 &  &  &  & 1\\
Dutch &  &  & 1 &  & 1\\
Museum & 1 &  & 1 & 1 & 0\\
Museumplein & 1 & 1 &  &  & 0\\
Van Gogh &  & 1 &  &  & 0\\
Stedelijk &  & 1 &  &  & 0\\
concertgebouw &  & 1 &  & 1 & 0\\
Rembrandt &  &  & 1 &  & 0\\
Vermeer &  &  & 1 &  & 0\\
Hals &  &  & 1 &  & 0
\end{tabular}

\item

$\vec{q}=(1,1,1,0,0,0,0,0,0,0)$

\[
\vec{doc1}\cdot\vec{q}=1+1+1+1+0+0+1+0+1+0+0+0=2
\]

\[
\vec{doc2}\cdot\vec{q}=1
\]

\[
\vec{doc3}\cdot\vec{q}=1
\]

\[
\vec{doc4}\cdot\vec{q}=1
\]

\item

\[
\cos(\widehat{doc1,q})=\frac{2}{\sqrt{4}\cdot\sqrt{3}}=0.5773...\approx0.577
\]

\[
\cos(\widehat{doc2,q})=\frac{1}{\sqrt{5}\cdot\sqrt{3}}=0.2581...\approx0.258
\]

\[
\cos(\widehat{doc3,q})=\frac{1}{\sqrt{5}\cdot\sqrt{3}}\approx0.258
\]

\[
\cos(\widehat{doc4,q})=\frac{1}{\sqrt{3}\cdot\sqrt{3}}=0.333...\approx0.333
\]

\item

For this query and document collection, the dot product similarity ranks doc1 as the most similar to the query, and doc 2,3,4 as equally similar to each other.

The cosine equality ranks most similar to be doc2=doc3$<$doc4$<$doc1

The difference in this case is not huge, both methods result in doc1 having the highest similarity. But the cosine similarity values rank 4 higher than 2,3, while 4 is intuitively less similar to the query.

\end{enumerate}

\item

\begin{enumerate}

\item
Doc1, doc2, doc3 are true positives, Doc4 is a false positive, it's about a concert, not the museum

\item
There are 5 relevant documents, 3 of which were retrieved, so there are $5-3=2$ false negatives.

There are 24 documents total, so $24-4-2=18$ true negatives

\item

\[
Precision=\frac{TP}{TP+FP}=\frac{3}{3+1}=\frac{3}{4}=0.75
\]

\[
Recall=\frac{TP}{TP+FN}=\frac{3}{3+2}=\frac{3}{5}=0.6
\]

\item

\[
F_{score}=2\cdot\frac{P\cdot R}{P+R}
\]

\[
=2\cdot\frac{0.75\cdot0.6}{0.75+0.6}
=0.666...\approx0.67
\]

\item

\[
idf_t=\log\left(\frac{N}{n_t}\right)
\]

For term Rijksmuseum:

\[
idf_t=\log\left(\frac{24}{10}\right)
\]

\[
tf-idf=1\cdot\log\left(\frac{24}{10}\right)=0.3802...\approx0.38
\]

For term Museumplein:

\[
idf_t=\log\left(\frac{24}{3}\right)
\]

\[
tf-idf=1\cdot\log\left(\frac{24}{3}\right)=0.9030...\approx0.90
\]

\item

Museumplein has a higher tf-idf value, this is due to the document frequency; museumplein occurs in less documents, so it gets weighted more.

\end{enumerate}

\item

\begin{enumerate}

\item

first, we compute the context from the sentence

\[
context=\{metal, collapsed, repeated, freezing, thawing, water\}
\]

Then we will compute the signature for every sense using the gloss and examples;

bridge$_1$: \{structure, people, vehicles, cross, obstacle, river, canal, railway, bridge, spans, Tacoma, Narrows, famously, collapsed, 1940\}

bridge$_2$: \{Upper, deck, ship, steered, captain, stands, wave, high, water, touched, bridge, modern, vessels, bridges, metal\}

Now for each sense, compute the overlap between context and the sense signature. The overlap count is the score.

bridge$_1$: overlap=1

bridge$_2$: overlap=2

\item

the sense bridge$_2$ would be the correct one, since the score from the simplified lesk algorithm is the highest.

\end{enumerate}

\end{enumerate}

\end{document}
